%% appendix_lit_review.tex — Comparison with Prior Habitability Frameworks
%% Moved from Introduction for word-count compliance.

\chapter{Comparison with Prior Habitability Frameworks}
\label{app:lit_review}

Several preceding studies have attempted to place planetary habitability on
a quantitative footing.  The \textit{Earth Similarity Index}
\citep{Schulze-Makuch2011} scores planetary bodies against Earth on
temperature, density, escape velocity, and surface chemistry, but
produces a single scalar per body and cannot represent spatial or temporal
variation.  \citet{Hendrix2019} reviewed the habitability potential of
icy ocean worlds using qualitative multi-criteria matrices, providing rich
scientific context but no spatial resolution or probabilistic framework.
Logistic regression and machine-learning habitability classifiers
(e.g.\ \citealt{Saha2018, Bora2016}) have been applied to exoplanet populations
but require labelled training data that does not exist for Titan.

The Bayesian approach adopted in this thesis differs from all of these: it
encodes scientific prior knowledge as probability distributions, updates them
with observational data via conjugate updating, and returns credible
intervals rather than point estimates: making uncertainty explicit
rather than implicit.  The choice of a Beta-conjugate (Beta-Binomial)
model is motivated by the bounded $[0,1]$ domain of habitability
probability and by the analytical tractability that permits simultaneous
evaluation at all 6.49 million grid pixels without numerical integration
or MCMC.  A Gaussian process or multivariate logistic regression could
in principle capture feature correlations more richly, but would require
either strong covariance priors or prohibitively dense labelled sampling
to estimate the correlation structure from Cassini data alone.  The
Beta-conjugate approach is preferred as the most tractable framework
that preserves physical interpretability at every pixel.

The closest methodological precedent is \citet{Affholder2021}, who applied
Bayesian statistical inference to Cassini INMS plume data from Enceladus,
comparing the observed methane and hydrogen escape rates against models of
abiotic serpentinisation and biotic methanogenesis.  Their finding (that
the observed methane escape rate is most parsimoniously explained by
biological methanogenesis if the prior probability of life is sufficient)
represents the most quantitatively rigorous Bayesian biosignature assessment
yet published for a solar system body.  The present work is directly inspired
by this approach, extended from a single-point plume observation to a full-globe,
multi-feature, multi-epoch spatial mapping.